% f2015hnw_q05_flow.tex — 2015 年高考数学·湖南卷·文科 第 5 题（程序框图）
%
% 【来源】`.work/phase3-b22/source/2015湖南文.pdf`（SHA-256
%   48dc0c6771ccea8291a3ac94336dd89a8b4ff89fa6baceaa76585c349582a074），
%   第 1 页（合订本页 957）；内容对象 5。逐路径实测见
%   `.work/phase3-b22/figure-inventory.md`。
%
% 【坐标口径】源 PDF 用户空间（原点左下、y 向上）。本图本地坐标
%   (x, y) = (x_page − 157.438, y_page − 122.447)，1 单位 = 1 TeX pt。
%   路径坐标逐点照录内容流（含三次贝塞尔控制点）。
%
% 元素计数：25 条路径（3 圆角框/2 平行四边框/3 矩形/1 菱形/箭头与回边）＋10 个文字标签
\begin{tikzpicture}[x=1pt, y=1pt,
    every node/.style={inner sep=0pt, font=\fontsize{10.5pt}{12.6pt}\selectfont}]
  \draw[line width=0.398pt] (70.2887,250.0260)   -- (50.2597,250.0260)   .. controls (48.0587,250.0260) and (46.2747,248.2420) .. (46.2747,246.0410)   -- (46.2747,237.1730)   .. controls (46.2747,234.9720) and (48.0587,233.1880) .. (50.2597,233.1880)   -- (70.2887,233.1880)   .. controls (72.4897,233.1880) and (74.2737,234.9720) .. (74.2737,237.1730)   -- (74.2737,246.0410)   .. controls (74.2737,248.2420) and (72.4897,250.0260) .. (70.2887,250.0260)   -- cycle;
  \draw[line width=0.398pt] (35.2437,201.7755)   -- (41.3717,218.6135)   -- (85.2987,218.6135)   -- (79.1707,201.7755)   -- cycle;
  \draw[line width=0.398pt] (28.9860,171.0010)   -- (91.5620,171.0010)   -- (91.5620,187.2020)   -- (28.9860,187.2020)   -- cycle;
  \draw[line width=0.398pt] (0.0000,125.7320)   -- (120.5470,125.7320)   -- (120.5470,156.4290)   -- (0.0000,156.4290)   -- cycle;
  \draw[line width=0.42pt] (45.6160,141.6850)   -- (115.8510,141.6850);
  \draw[line width=0.398pt] (37.1150,96.3610)   -- (83.4330,96.3610)   -- (83.4330,111.1600)   -- (37.1150,111.1600)   -- cycle;
  \draw[line width=0.398pt] (97.7278,72.4052)   -- (60.2738,81.7072)   -- (22.8198,72.4052)   -- (60.2738,63.1032)   -- cycle;
  \draw[line width=0.398pt] (34.8656,31.5885)   -- (41.0016,48.4475)   -- (85.6766,48.4475)   -- (79.5416,31.5885)   -- cycle;
  \draw[line width=0.398pt] (70.2887,17.0158)   -- (50.2597,17.0158)   .. controls (48.0587,17.0158) and (46.2747,15.2318) .. (46.2747,13.0308)   -- (46.2747,3.9848)   .. controls (46.2747,1.7838) and (48.0587,-0.0002) .. (50.2597,-0.0002)   -- (70.2887,-0.0002)   .. controls (72.4897,-0.0002) and (74.2737,1.7838) .. (74.2737,3.9848)   -- (74.2737,13.0308)   .. controls (74.2737,15.2318) and (72.4897,17.0158) .. (70.2887,17.0158)   -- cycle;
  \draw[line width=0.398pt] (60.2739,232.9888)   -- (60.2739,222.4008);
  \fill[line width=1.0pt] (60.2736,218.8142)   .. controls (60.4726,219.8772) and (61.0706,221.6042) .. (61.7676,222.7992)   -- (58.7796,222.7992)   .. controls (59.4766,221.6042) and (60.0746,219.8772) .. (60.2736,218.8142);
  \draw[line width=0.398pt] (60.2740,201.5755)   -- (60.2730,190.9885);
  \fill[line width=1.0pt] (60.2733,187.4017)   .. controls (60.4723,188.4646) and (61.0703,190.1916) .. (61.7674,191.3866)   -- (58.7794,191.3867)   .. controls (59.4763,190.1917) and (60.0743,188.4647) .. (60.2733,187.4017);
  \draw[line width=0.398pt] (60.2740,170.8022)   -- (60.2740,160.2152);
  \fill[line width=1.0pt] (60.2740,156.6288)   .. controls (60.4730,157.6918) and (61.0710,159.4188) .. (61.7680,160.6138)   -- (58.7800,160.6138)   .. controls (59.4770,159.4188) and (60.0750,157.6918) .. (60.2740,156.6288);
  \draw[line width=0.398pt] (60.2740,125.5332)   -- (60.2740,114.9462);
  \fill[line width=1.0pt] (60.2740,111.3599)   .. controls (60.4730,112.4229) and (61.0710,114.1499) .. (61.7680,115.3449)   -- (58.7800,115.3449)   .. controls (59.4770,114.1499) and (60.0750,112.4229) .. (60.2740,111.3599);
  \draw[line width=0.398pt] (60.2740,96.1619)   -- (60.2740,85.5749);
  \fill[line width=1.0pt] (60.2740,81.9880)   .. controls (60.4730,83.0510) and (61.0710,84.7780) .. (61.7680,85.9730)   -- (58.7800,85.9730)   .. controls (59.4770,84.7780) and (60.0750,83.0510) .. (60.2740,81.9880);
  \draw[line width=0.398pt] (60.2740,62.8220)   -- (60.2740,52.2350);
  \fill[line width=1.0pt] (60.2743,48.6487)   .. controls (60.4733,49.7117) and (61.0713,51.4387) .. (61.7683,52.6337)   -- (58.7803,52.6337)   .. controls (59.4773,51.4387) and (60.0753,49.7117) .. (60.2743,48.6487);
  \draw[line width=0.398pt] (60.2750,31.3882)   -- (60.2740,20.8012);
  \fill[line width=1.0pt] (60.2734,17.2151)   .. controls (60.4725,18.2781) and (61.0706,20.0050) .. (61.7677,21.2000)   -- (58.7797,21.2002)   .. controls (59.4766,20.0052) and (60.0745,18.2781) .. (60.2734,17.2151);
  \draw[line width=0.398pt] (98.0095,72.4052)   -- (132.0255,72.4052)   -- (132.0255,163.7152)   -- (63.8605,163.7152);
  \fill[line width=1.0pt] (60.2740,163.7155)   .. controls (61.3370,163.5165) and (63.0640,162.9185) .. (64.2590,162.2215)   -- (64.2590,165.2095)   .. controls (63.0640,164.5125) and (61.3370,163.9145) .. (60.2740,163.7155);

  \node[anchor=center] at (60.272,242.193) {开始};  % 源中心 (217.71,364.64)
  \node[anchor=center] at (60.262,210.253) {输入 $n$};  % 源中心 (217.70,332.70)
  \node[anchor=center] at (60.262,179.153) {$i=1,\ S=0$};  % 源中心 (217.70,301.60)
  \node[anchor=center] at (60.262,141.053) {$S=S+\dfrac{1}{(2i-1)(2i+1)}$};  % 源中心 (217.70,263.50)
  \node[anchor=center] at (60.312,103.363) {$i=i+1$};  % 源中心 (217.75,225.81)
  \node[anchor=center] at (60.262,71.533) {$i>n?$};  % 源中心 (217.70,193.98)
  \node[anchor=center] at (59.962,40.053) {输出 $S$};  % 源中心 (217.40,162.50)
  \node[anchor=center] at (60.272,9.113) {结束};  % 源中心 (217.71,131.56)
  \node[anchor=center] at (69.222,56.423) {是};  % 源中心 (226.66,178.87)
  \node[anchor=center] at (98.012,81.373) {否};  % 源中心 (255.45,203.82)

\end{tikzpicture}
%
% 【自验结果】
%   探针：latexmk -pdfxe（待跑）。
